\section{Conformal radius and the law centered at one point}
\label{op:section}

Our first goal is a finite change of law toward one interior point.
The same calculation will supply a conditional Green potential, uniform
angle mixing, and a finite physical lifetime. These are the inputs for
both the mass construction and the eventual tests at successive radii.

Fix $0<\kappa<8$ throughout, and put
\begin{equation}
 a=\frac2\kappa>\frac14,\qquad d=1+\frac1{4a},\qquad
 \alpha=2-d\in(0,1),\qquad p=4a-1>0.
 \label{op:parameters}
\end{equation}
In particular, $2a\alpha=p/2$. If $t_{\rm cap}$ denotes the
capacity-$2t_{\rm cap}$ time of Theorem~\ref{thm:main}, put
$u=\kappa t_{\rm cap}$ and
$B^{\rm int}_u=\sqrt\kappa B^{\rm cap}_{u/\kappa}$.
This driver is standard Brownian motion. Writing $t$ for the internal
time $u$ and $B$ for $B^{\rm int}$ below, the equation becomes
\begin{equation}
 \partial_t g_t(z)=\frac{a}{g_t(z)-B_t},\qquad g_0(z)=z,
 \label{op:loewner}
\end{equation}
with a standard Brownian driver and half-plane capacity $at$.
The original normalization and continuous trace generation are given by
Rohde--Schramm \cite[Section~2.1 and Theorem~5.1]{RohdeSchramm2005}.
A constant change of time preserves the
assertion about all nontrivial finite time intervals.

We first specify the geometric and probabilistic conventions. Write $H_t$
for the survival domain of \eqref{op:loewner}. It is the unbounded component
of $\HH\setminus\gamma([0,t])$, and $g_t:H_t\to\HH$ is conformal.
Moreover,
\begin{equation}
 \gamma([t,\infty))\subset\overline{H_t}.
 \label{op:future-trace}
\end{equation}
These are the filled-hull conventions for a Loewner evolution generated
by a continuous trace. The domain identification is part of
Rohde--Schramm \cite[Theorem~5.1]{RohdeSchramm2005}, including when
the trace has self-touches. To obtain \eqref{op:future-trace}, for
$u\ge t$ and $\varepsilon>0$ observe that
$g_u^{-1}(B_u+i\varepsilon)\in H_u\subset H_t$.
The inverse-map definition of $\gamma(u)$ and the limit
$\varepsilon\downarrow0$ give the asserted containment.

We also work in a fixed simply connected proper domain
$D=J(\HH)$, where $J:\HH\to D$ is conformal, and set $D_t=J(H_t)$.
Only the interior domain evolution is needed for an arbitrary $J$;
no boundary extension of $J$ is being assumed. At a finite stopping time
$S$, the map
\[
 J_S=J\circ(g_S-B_S)^{-1}:\HH\longrightarrow D_S
\]
identifies the remaining domain with a fresh copy of the half-plane.
Continuing the ODE gives the composition rule, and the strong Markov
property gives the fresh driver $B_{S+t}-B_S$.

For a simply connected domain $U$ and $x\in U$, our conformal radius is
$R_U(x)=|F'(0)|$, where $F:\DD\to U$ is conformal with $F(0)=x$.
We use the classical Koebe estimates in this convention
\cite[Theorem~3.2.5, Corollary~3.2.6, and Theorems~3.2.9
and~3.2.11]{Beliaev2015Notes}:
\begin{align}
 R_U(x)/4&\le\operatorname{dist}(x,\partial U)\le R_U(x),
 \label{op:koebe-distance}\\
 R_U(x)\frac{r}{(1+r)^2}
 &\le |F(w)-x|\le R_U(x)\frac{r}{(1-r)^2},
 \quad |w|=r<1, \label{op:koebe-growth}\\
 R_U(x)\frac{1-r}{(1+r)^3}
 &\le |F'(w)|\le R_U(x)\frac{1+r}{(1-r)^3},
 \quad |w|\le r<1. \label{op:koebe-derivative}
\end{align}

All path functionals used under a changed law are taken on the canonical
continuous-driver space with its \emph{raw} natural filtration
$(\mathcal F_t)$. Usual augmentations may be used for stochastic-calculus
arguments under the original Wiener law, followed by raw versions of the
resulting adapted quantities. The centered laws below are laws of finite
stopped data. Their consistent limits do not assert absolute continuity
on the original completed infinite-future sigma field. In particular,
later adapted mass functionals will be evaluated using one raw
nonanticipating version, rather than choosing separate versions for
different centers.

\subsection{The local martingale and finite stopped laws}

The finite-radius Green tilt and its centered angular diffusion are
standard in the study of two-sided radial SLE; see
Lawler--Zhou \cite[Section~2.1]{LawlerZhou2013} and
Lawler--Rezaei \cite[Section~4.2, equations~(38)--(40)]{LawlerRezaei2015}.
We give the cutoff argument and uniform mixing proof in the present
full-conformal-radius normalization, with all changes of law made at
finite stopped levels.

Fix $x\in D$, put $z=J^{-1}(x)$, and let $T_x$ be its lifetime. Until
$T_x$, write
\[
 Z_t=X_t+iY_t=g_t(z)-B_t,\qquad
 \theta_t=\arg Z_t,\qquad V_t=\frac{Y_t}{|Z_t|^2},\qquad
 s_x(t)=\int_0^t V_u^2\,\dd u.
\]
Set $R_t(x)=R_{D_t}(x)$. Conformal covariance of radius gives
\[
 R_t(x)=\frac{2Y_t|J'(z)|}{|g_t'(z)|}.
\]
Since $\dd Z_t=aZ_t^{-1}\dd t-\dd B_t$, differentiation of this
formula and It\^o's Formula for $\log Z_t$ give
\begin{align}
 \dd\log R_t(x)&=-2aV_t^2\,\dd t,
 &R_t(x)&=R_0(x)e^{-2as_x(t)},\label{op:radius-clock}\\
 \dd\theta_t&=V_t\,\dd B_t+(1-2a)\cot\theta_t\,V_t^2\,\dd t.
 \label{op:angle-sde}
\end{align}
Thus the angle in inner time has generator
\[
 \mathcal L=\frac12\partial_\theta^2
                +(1-2a)\cot\theta\,\partial_\theta.
\]
For every real $r$, direct differentiation yields
\begin{equation}
 \frac{\mathcal L\sin^r\theta}{\sin^r\theta}
 =-\frac r2+\frac{r(r-p)}2\cot^2\theta.
 \label{op:sine-generator}
\end{equation}
Consequently
\begin{equation}
 G_t(x)=R_t(x)^{-\alpha}\sin^p\theta_t(x)
 \quad (x\in D_t),\qquad G_t(x)=0\quad(x\notin D_t)
 \label{op:green}
\end{equation}
is a local martingale during survival, with stochastic differential
$\dd G_t=G_t p\cot\theta_t V_t\,\dd B_t$.

For $s\ge0$, let $\tau_s(x)$ be the time at which the radius first
attains $R_0(x)e^{-2as}$ while $x$ is alive; it is infinity if this
never occurs. This is the inverse inner clock on successful paths.

\begin{lemma}[Finite stopped change of law]\label{op:finite-tilt}
For every finite $s\ge0$, the nonnegative weight
\begin{equation}
 M_s(x)=\frac{G_{\tau_s(x)}(x)}{G_0(x)}
                 \ind_{\{\tau_s(x)<\infty\}}
 \label{op:finite-density}
\end{equation}
has expectation one. It defines a law $\PP_x^*$ on the path stopped at
$\tau_s(x)$, consistently as $s$ increases. Under these laws, the angle
in inner time solves
\begin{equation}
 \dd\theta_s=\dd W_s+2a\cot\theta_s\,\dd s,
 \qquad
 \mathcal L^*=\tfrac12\partial_\theta^2
                       +2a\cot\theta\,\partial_\theta,
 \label{op:tilted-angle}
\end{equation}
and neither endpoint is reached in finite inner time.

For an arbitrary stopping time $S$, on $\{S<\infty\}$ one has
\begin{equation}
 \EE\!\left[G_{\tau_s(x)}(x)
       \ind_{\{\tau_s(x)<\infty,\ S<\tau_s(x)\}}
                  \mid\mathcal F_S\right]
 =G_S(x)\ind_{\{x\in D_S,\ s_x(S)<s\}}.
 \label{op:conditional-weight}
\end{equation}
Conditionally on such a surviving state, the remainder of the centered
law is the finite stopped change of law in the domain $D_S$, with
remaining inner duration $s-s_x(S)$. Indeed,
\[
 R_S(x)e^{-2a(s-s_x(S))}=R_0(x)e^{-2as},
\]
so the restarted target is the original radius level.
\end{lemma}

\begin{proof}
Localize first by stopping when the angle exits $(1/N,\pi-1/N)$,
when $Y$ falls below $1/N$, or when physical time reaches $N$.
Write $\zeta_N$ for this cutoff and $\PP_x^{*,N}$ for its changed law.
The stochastic coefficient of $G/G_0$ is bounded on this stopped
interval. Girsanov and \eqref{op:angle-sde} change the inner-time drift
from $(1-2a)\cot\theta$ to $(1-2a+p)\cot\theta=2a\cot\theta$.

We verify that removing the cutoffs loses no probability at any finite
inner horizon. For $0<q<p$, set $F_q(\theta)=\sin^{-q}\theta$.
The new generator satisfies
\begin{equation}
 \mathcal L^*F_q
 =F_q\left(\frac q2-\frac{q(p-q)}2\cot^2\theta\right)
 \le\frac q2 F_q.
 \label{op:lyapunov}
\end{equation}
The locally defined diffusion \eqref{op:tilted-angle} therefore obeys,
for its first exit $\sigma_\delta$ from $(\delta,\pi-\delta)$,
\[
 \PP^*_{\theta_0}(\sigma_\delta\le s)
 \le e^{qs/2}F_q(\theta_0)\sin^q\delta.
\]
This follows by stopping the nonnegative local supermartingale
$e^{-qu/2}F_q(\theta_u)$. The bound tends to zero and proves global
existence in inner time and endpoint nonattainment.

The physical clock cannot obstruct this argument. Along every compact
inner interval on which the angle stays in $(0,\pi)$,
\begin{equation}
 \frac{\dd}{\dd s}\log Y_s=-a\sin^{-2}\theta_s,
 \qquad
 \frac{\dd t}{\dd s}=Y_s^2\sin^{-4}\theta_s.
 \label{op:physical-clock}
\end{equation}
The first equation follows from $\dd Y_t=-aY_t|Z_t|^{-2}\dd t$;
the second is $V_t^{-2}$, using $V_t=\sin^2\theta_t/Y_t$.
These formulas keep $Y$ strictly positive and physical time finite on
that compact interval. Thus, for $E_N=\{\tau_s<\zeta_N\}$, the
localized changed laws satisfy $\PP_x^{*,N}(E_N)\to1$.
On $E_N$ the localized density at the target is exactly $M_s$, so
\[
 \EE[M_s\ind_{E_N}]=\PP_x^{*,N}(E_N)\longrightarrow1.
\]
Under the original law, $E_N$ increases to $\{\tau_s<\infty\}$:
a point alive through its finite target time has positive height and
angle minima on that compact interval. Monotone convergence therefore
gives $\EE M_s=1$. This calculation uses the localized density only
before the cutoff; it does not continue the weight through a lifetime.

Repeating this calculation after $S$, using the Loewner composition
rule and strong Markov property, gives \eqref{op:conditional-weight}.
If the point is already dead, or the level has already been reached,
both sides are zero. Applying the same calculation between two finite
inner levels proves consistency and the asserted conditional restart.
Only stopped paths through a finite inner level enter any of these
changes of law.
\end{proof}

\subsection{Uniform mixing and the remaining potential}

The stationary density and uniform mixing are also described in
\cite[Section~4.2]{LawlerRezaei2015}. The orthogonal-polynomial
proof has a direct counterpart in
Zhan \cite[Proposition~4.2]{Zhan2019Optimal}. Writing $X$ for the
diffusion in that proposition, the substitution $\theta_s=X_{4s/\kappa}/2$
gives our centered diffusion. We prove mixing uniformly over the initial
angle and give the finite-lifetime estimates needed for the subsequent
potential identity.

\begin{lemma}[Uniform angle mixing]\label{op:mixing}
Let
\begin{equation}
 \pi_a(\dd\theta)=c_a\sin^{4a}\theta\,\dd\theta,
 \qquad c_a^{-1}=\int_0^\pi\sin^{4a}\theta\,\dd\theta.
 \label{op:stationary}
\end{equation}
There is a deterministic $s_a<\infty$ such that, for all $s\ge s_a$
and every initial angle in $(0,\pi)$, the transition law of
\eqref{op:tilted-angle} has density between $1/2$ and $2$ relative to
$\pi_a$. In particular, after such a gap the conditional expectation
of any nonnegative angle function lies between one half and twice its
$\pi_a$ integral. Moreover,
\begin{equation}
 \int_0^\pi\sin^{-p}\theta\,\pi_a(\dd\theta)
 =c_a\int_0^\pi\sin\theta\,\dd\theta<\infty.
 \label{op:inverse-angle}
\end{equation}
\end{lemma}

\begin{proof}
For $v=\cos\theta$, the generator and candidate invariant probability
measure become
\[
 \mathcal A=\frac{1-v^2}{2}\partial_v^2
              -\left(\frac12+2a\right)v\partial_v,
 \qquad \varpi_a(\dd v)=\widetilde c_a
              (1-v^2)^{2a-1/2}\,\dd v.
\]
Integration by parts makes $\mathcal A$ symmetric on polynomials; the
boundary term vanishes because $(1-v^2)^{2a+1/2}$ tends to zero.
Let $P_m$ be the orthonormal polynomials for $\varpi_a$, with $P_0=1$.
The operator preserves degree and acts on the leading degree-$m$
coefficient by $-\lambda_m$, where
$\lambda_m=m(m+4a)/2$. Symmetry and orthogonality therefore give
$\mathcal AP_m=-\lambda_mP_m$.

Here is a sufficient elementary bound on these polynomials. Choose
$m+1$ disjoint intervals of length comparable to $1/(m+1)$, separated
by the same order, in $[-1/2,1/2]$. The density of $\varpi_a$ is bounded
below there. The unit $L^2$ norm supplies in each interval a point
where $|P_m|\le C_a\sqrt{m+1}$. Lagrange interpolation at these points,
whose pairwise separation is at least $c/(m+1)$, gives
\[
 \|P_m\|_{\infty,[-1,1]}
 \le \exp\{C_a(m+1)\log(m+2)\}.
\]
Consequently, for every $s>0$ the series
\[
 k_s(v,w)=1+\sum_{m\ge1}e^{-s\lambda_m}P_m(v)P_m(w)
\]
converges absolutely and uniformly in $(v,w)\in[-1,1]^2$.
It\^o's Formula for bounded polynomials gives
$\EE_vP_m(v_s)=e^{-s\lambda_m}P_m(v)$, so the actual transition
probability and $k_s(v,w)\varpi_a(\dd w)$ have identical polynomial
moments. Polynomial density in $C([-1,1])$ identifies the measures;
this also proves that the displayed kernel represents a nonnegative
measure. Its nonconstant part tends uniformly to zero as
$s\to\infty$, proving the bounds. The conditional assertion follows
from the Markov property. Finally $4a-p=1$ proves
\eqref{op:inverse-angle}.
\end{proof}

\begin{proposition}[Finite lifetime and tested potential identity]
\label{op:potential}
Under the consistent centered laws,
$T^*(x)=\lim_{s\to\infty}\tau_s(x)$ is finite almost surely.
For every finite deterministic $t$ and bounded raw
$\mathcal F_t$-measurable random variable $H$,
\begin{equation}
 \EE[H G_t(x)]
 =G_0(x)\EE_x^*[H\ind_{\{t<T^*(x)\}}].
 \label{op:tested-potential}
\end{equation}
On the right, $H$ is evaluated only on stopped data extending beyond
$t$. Thus
\begin{equation}
 \EE G_t(x)=G_0(x)\PP_x^*(T^*(x)>t)\longrightarrow0
 \quad(t\to\infty).
 \label{op:potential-decay}
\end{equation}
For every stopping time $S$,
\begin{equation}
 \EE[G_S(x)\ind_{\{S<\infty\}}]\le G_0(x).
 \label{op:optional-bound}
\end{equation}
\end{proposition}

\begin{proof}
Choose $0<q<p$ and $\varepsilon\in(0,a/2)$. Lemma~\ref{op:mixing}
gives a uniform bound on $\EE_x^*F_q(\theta_m)$ for all sufficiently
large integers $m$, since $F_q$ is integrable against $\pi_a$.
The nonnegative supermartingale from \eqref{op:lyapunov} and its
maximal inequality imply
\[
 \PP_x^*\!\left(\sup_{m\le s\le m+1}F_q(\theta_s)
                      >e^{q\varepsilon m}\right)
 \le C e^{-q\varepsilon m}.
\]
Borel--Cantelli shows that eventually
$\sin^{-4}\theta_s\le e^{4\varepsilon s}$.
Equation~\eqref{op:physical-clock} also gives $Y_s\le Y_0e^{-as}$.
Therefore
\[
 T^*(x)=\int_0^\infty Y_s^2\sin^{-4}\theta_s\,\dd s<\infty.
\]
The integral over each initial compact interval is finite by endpoint
nonattainment; the remaining integrand is bounded by a constant times
$e^{-(2a-4\varepsilon)s}$.

For a finite inner horizon $l$, Lemma~\ref{op:finite-tilt} gives
\[
 G_0(x)\EE_x^*[H\ind_{\{t<\tau_l(x)\}}]
 =\EE[H G_t(x)\ind_{\{x\in D_t,\ s_x(t)<l\}}].
\]
First take $H=1$. Monotone convergence as $l\to\infty$ proves the
untested identity and the integrable bound $\EE G_t\le G_0$.
The same passage for bounded signed $H$ then follows by dominated
convergence. The union of the events on the left is exactly
$\{t<T^*\}$; no value at $T^*$ is used. Decay follows from
$T^*<\infty$. Applying \eqref{op:conditional-weight} at $S$, taking
expectations, and increasing $l$ proves \eqref{op:optional-bound}.

These identities also make the parameter measurability needed below
explicit. For a raw $A\in\mathcal F_t$,
\[
 \PP_x^*(A,t<T^*)
 =\lim_{l\to\infty}\frac1{G_0(x)}
   \EE[\ind_A G_t(x)
                 \ind_{\{x\in D_t,\ s_x(t)<l\}}].
\]
The ODE, its spatial derivative, and the inverse clock are measurable
in the driver and the marked point. The right side is consequently a
measurable kernel in $x$. This permits integration over centers and
does not require a common exceptional-set assertion for all centers.
\end{proof}
